\IfFileExists{stacks-project.cls}{%
\documentclass{stacks-project}
}{%
\documentclass{amsart}
}

% For dealing with references we use the comment environment
\usepackage{verbatim}
\newenvironment{reference}{\comment}{\endcomment}
%\newenvironment{reference}{}{}
\newenvironment{slogan}{\comment}{\endcomment}
\newenvironment{history}{\comment}{\endcomment}

% For commutative diagrams we use Xy-pic
\usepackage[all]{xy}

% We use 2cell for 2-commutative diagrams.
\xyoption{2cell}
\UseAllTwocells

% We use multicol for the list of chapters between chapters
\usepackage{multicol}

% This is generally recommended for better output
\usepackage{lmodern}
\usepackage[T1]{fontenc}

% For cross-file-references
\usepackage{xr-hyper}

% Package for hypertext links:
\usepackage{hyperref}

% For any local file, say "hello.tex" you want to link to please
% use \externaldocument[hello-]{hello}
\externaldocument[introduction-]{introduction}
\externaldocument[conventions-]{conventions}
\externaldocument[sets-]{sets}
\externaldocument[categories-]{categories}
\externaldocument[topology-]{topology}
\externaldocument[sheaves-]{sheaves}
\externaldocument[sites-]{sites}
\externaldocument[stacks-]{stacks}
\externaldocument[fields-]{fields}
\externaldocument[algebra-]{algebra}
\externaldocument[brauer-]{brauer}
\externaldocument[homology-]{homology}
\externaldocument[derived-]{derived}
\externaldocument[simplicial-]{simplicial}
\externaldocument[more-algebra-]{more-algebra}
\externaldocument[smoothing-]{smoothing}
\externaldocument[modules-]{modules}
\externaldocument[sites-modules-]{sites-modules}
\externaldocument[injectives-]{injectives}
\externaldocument[cohomology-]{cohomology}
\externaldocument[sites-cohomology-]{sites-cohomology}
\externaldocument[dga-]{dga}
\externaldocument[dpa-]{dpa}
\externaldocument[sdga-]{sdga}
\externaldocument[hypercovering-]{hypercovering}
\externaldocument[schemes-]{schemes}
\externaldocument[constructions-]{constructions}
\externaldocument[properties-]{properties}
\externaldocument[morphisms-]{morphisms}
\externaldocument[coherent-]{coherent}
\externaldocument[divisors-]{divisors}
\externaldocument[limits-]{limits}
\externaldocument[varieties-]{varieties}
\externaldocument[topologies-]{topologies}
\externaldocument[descent-]{descent}
\externaldocument[perfect-]{perfect}
\externaldocument[more-morphisms-]{more-morphisms}
\externaldocument[flat-]{flat}
\externaldocument[groupoids-]{groupoids}
\externaldocument[more-groupoids-]{more-groupoids}
\externaldocument[etale-]{etale}
\externaldocument[chow-]{chow}
\externaldocument[intersection-]{intersection}
\externaldocument[pic-]{pic}
\externaldocument[weil-]{weil}
\externaldocument[adequate-]{adequate}
\externaldocument[dualizing-]{dualizing}
\externaldocument[duality-]{duality}
\externaldocument[discriminant-]{discriminant}
\externaldocument[derham-]{derham}
\externaldocument[local-cohomology-]{local-cohomology}
\externaldocument[algebraization-]{algebraization}
\externaldocument[curves-]{curves}
\externaldocument[resolve-]{resolve}
\externaldocument[models-]{models}
\externaldocument[functors-]{functors}
\externaldocument[equiv-]{equiv}
\externaldocument[pione-]{pione}
\externaldocument[etale-cohomology-]{etale-cohomology}
\externaldocument[proetale-]{proetale}
\externaldocument[relative-cycles-]{relative-cycles}
\externaldocument[more-etale-]{more-etale}
\externaldocument[trace-]{trace}
\externaldocument[crystalline-]{crystalline}
\externaldocument[spaces-]{spaces}
\externaldocument[spaces-properties-]{spaces-properties}
\externaldocument[spaces-morphisms-]{spaces-morphisms}
\externaldocument[decent-spaces-]{decent-spaces}
\externaldocument[spaces-cohomology-]{spaces-cohomology}
\externaldocument[spaces-limits-]{spaces-limits}
\externaldocument[spaces-divisors-]{spaces-divisors}
\externaldocument[spaces-over-fields-]{spaces-over-fields}
\externaldocument[spaces-topologies-]{spaces-topologies}
\externaldocument[spaces-descent-]{spaces-descent}
\externaldocument[spaces-perfect-]{spaces-perfect}
\externaldocument[spaces-more-morphisms-]{spaces-more-morphisms}
\externaldocument[spaces-flat-]{spaces-flat}
\externaldocument[spaces-groupoids-]{spaces-groupoids}
\externaldocument[spaces-more-groupoids-]{spaces-more-groupoids}
\externaldocument[bootstrap-]{bootstrap}
\externaldocument[spaces-pushouts-]{spaces-pushouts}
\externaldocument[spaces-chow-]{spaces-chow}
\externaldocument[groupoids-quotients-]{groupoids-quotients}
\externaldocument[spaces-more-cohomology-]{spaces-more-cohomology}
\externaldocument[spaces-simplicial-]{spaces-simplicial}
\externaldocument[spaces-duality-]{spaces-duality}
\externaldocument[formal-spaces-]{formal-spaces}
\externaldocument[restricted-]{restricted}
\externaldocument[spaces-resolve-]{spaces-resolve}
\externaldocument[formal-defos-]{formal-defos}
\externaldocument[defos-]{defos}
\externaldocument[cotangent-]{cotangent}
\externaldocument[examples-defos-]{examples-defos}
\externaldocument[algebraic-]{algebraic}
\externaldocument[examples-stacks-]{examples-stacks}
\externaldocument[stacks-sheaves-]{stacks-sheaves}
\externaldocument[criteria-]{criteria}
\externaldocument[artin-]{artin}
\externaldocument[quot-]{quot}
\externaldocument[stacks-properties-]{stacks-properties}
\externaldocument[stacks-morphisms-]{stacks-morphisms}
\externaldocument[stacks-limits-]{stacks-limits}
\externaldocument[stacks-cohomology-]{stacks-cohomology}
\externaldocument[stacks-perfect-]{stacks-perfect}
\externaldocument[stacks-introduction-]{stacks-introduction}
\externaldocument[stacks-more-morphisms-]{stacks-more-morphisms}
\externaldocument[stacks-geometry-]{stacks-geometry}
\externaldocument[moduli-]{moduli}
\externaldocument[moduli-curves-]{moduli-curves}
\externaldocument[examples-]{examples}
\externaldocument[exercises-]{exercises}
\externaldocument[guide-]{guide}
\externaldocument[desirables-]{desirables}
\externaldocument[coding-]{coding}
\externaldocument[obsolete-]{obsolete}
\externaldocument[fdl-]{fdl}
\externaldocument[index-]{index}

% Theorem environments.
%
\theoremstyle{plain}
\newtheorem{theorem}[subsection]{Theorem}
\newtheorem{proposition}[subsection]{Proposition}
\newtheorem{lemma}[subsection]{Lemma}

\theoremstyle{definition}
\newtheorem{definition}[subsection]{Definition}
\newtheorem{example}[subsection]{Example}
\newtheorem{exercise}[subsection]{Exercise}
\newtheorem{situation}[subsection]{Situation}

\theoremstyle{remark}
\newtheorem{remark}[subsection]{Remark}
\newtheorem{remarks}[subsection]{Remarks}

\numberwithin{equation}{subsection}

% Macros
%
\def\lim{\mathop{\mathrm{lim}}\nolimits}
\def\colim{\mathop{\mathrm{colim}}\nolimits}
\def\Spec{\mathop{\mathrm{Spec}}}
\def\Hom{\mathop{\mathrm{Hom}}\nolimits}
\def\Ext{\mathop{\mathrm{Ext}}\nolimits}
\def\SheafHom{\mathop{\mathcal{H}\!\mathit{om}}\nolimits}
\def\SheafExt{\mathop{\mathcal{E}\!\mathit{xt}}\nolimits}
\def\Sch{\mathit{Sch}}
\def\Mor{\mathop{\mathrm{Mor}}\nolimits}
\def\Ob{\mathop{\mathrm{Ob}}\nolimits}
\def\Sh{\mathop{\mathit{Sh}}\nolimits}
\def\NL{\mathop{N\!L}\nolimits}
\def\CH{\mathop{\mathrm{CH}}\nolimits}
\def\proetale{{pro\text{-}\acute{e}tale}}
\def\etale{{\acute{e}tale}}
\def\QCoh{\mathit{QCoh}}
\def\Ker{\mathop{\mathrm{Ker}}}
\def\Im{\mathop{\mathrm{Im}}}
\def\Coker{\mathop{\mathrm{Coker}}}
\def\Coim{\mathop{\mathrm{Coim}}}

% Boxtimes
%
\DeclareMathSymbol{\boxtimes}{\mathbin}{AMSa}{"02}

%
% Macros for moduli stacks/spaces
%
\def\QCohstack{\mathcal{QC}\!\mathit{oh}}
\def\Cohstack{\mathcal{C}\!\mathit{oh}}
\def\Spacesstack{\mathcal{S}\!\mathit{paces}}
\def\Quotfunctor{\mathrm{Quot}}
\def\Hilbfunctor{\mathrm{Hilb}}
\def\Curvesstack{\mathcal{C}\!\mathit{urves}}
\def\Polarizedstack{\mathcal{P}\!\mathit{olarized}}
\def\Complexesstack{\mathcal{C}\!\mathit{omplexes}}
% \Pic is the operator that assigns to X its picard group, usage \Pic(X)
% \Picardstack_{X/B} denotes the Picard stack of X over B
% \Picardfunctor_{X/B} denotes the Picard functor of X over B
\def\Pic{\mathop{\mathrm{Pic}}\nolimits}
\def\Picardstack{\mathcal{P}\!\mathit{ic}}
\def\Picardfunctor{\mathrm{Pic}}
\def\Deformationcategory{\mathcal{D}\!\mathit{ef}}
\usepackage{cleveref}
\crefname{originallemma}{original lemma}{original lemmas}
\Crefname{originallemma}{Original lemma}{Original lemmas}
\makeatletter
\newlabel{lemma-flat-proper-perfect-direct-image-general}{{0B91}{source}{}{sourcepointer-1}{}}
\newlabel{lemma-flat-proper-perfect-direct-image-general@cref}{{[originallemma][0][]0B91}{[1][0][]source}{}{}{}}
\makeatother
\crefname{originallemma}{original lemma}{original lemmas}
\Crefname{originallemma}{Original lemma}{Original lemmas}
\makeatletter
\newlabel{coherent-lemma-flat-base-change-cohomology}{{02KH}{source}{}{sourcepointer-2}{}}
\newlabel{coherent-lemma-flat-base-change-cohomology@cref}{{[originallemma][0][]02KH}{[1][0][]source}{}{}{}}
\makeatother
\crefname{originallemma}{original lemma}{original lemmas}
\Crefname{originallemma}{Original lemma}{Original lemmas}
\makeatletter
\newlabel{more-algebra-lemma-better-cut-complex-in-two}{{0A1U}{source}{}{sourcepointer-3}{}}
\newlabel{more-algebra-lemma-better-cut-complex-in-two@cref}{{[originallemma][0][]0A1U}{[1][0][]source}{}{}{}}
\makeatother
\crefname{originallemma}{original lemma}{original lemmas}
\Crefname{originallemma}{Original lemma}{Original lemmas}
\makeatletter
\newlabel{more-algebra-lemma-perfect}{{0658}{source}{}{sourcepointer-4}{}}
\newlabel{more-algebra-lemma-perfect@cref}{{[originallemma][0][]0658}{[1][0][]source}{}{}{}}
\makeatother
\crefname{originallemma}{original lemma}{original lemmas}
\Crefname{originallemma}{Original lemma}{Original lemmas}
\makeatletter
\newlabel{varieties-lemma-proper-geometrically-reduced-global-sections}{{0BUG}{source}{}{sourcepointer-5}{}}
\newlabel{varieties-lemma-proper-geometrically-reduced-global-sections@cref}{{[originallemma][0][]0BUG}{[1][0][]source}{}{}{}}
\makeatother
\crefname{originallemma}{original lemma}{original lemmas}
\Crefname{originallemma}{Original lemma}{Original lemmas}
\makeatletter
\newlabel{lemma-proper-idempotent-on-fibre}{{0G7X}{source}{}{sourcepointer-6}{}}
\newlabel{lemma-proper-idempotent-on-fibre@cref}{{[originallemma][0][]0G7X}{[1][0][]source}{}{}{}}
\makeatother
\crefname{originallemma}{original lemma}{original lemmas}
\Crefname{originallemma}{Original lemma}{Original lemmas}
\makeatletter
\newlabel{morphisms-lemma-set-points-where-fibres-etale}{{0476}{source}{}{sourcepointer-7}{}}
\newlabel{morphisms-lemma-set-points-where-fibres-etale@cref}{{[originallemma][0][]0476}{[1][0][]source}{}{}{}}
\makeatother
\crefname{originallemma}{original lemma}{original lemmas}
\Crefname{originallemma}{Original lemma}{Original lemmas}
\makeatletter
\newlabel{algebra-lemma-finite-projective}{{00NX}{source}{}{sourcepointer-8}{}}
\newlabel{algebra-lemma-finite-projective@cref}{{[originallemma][0][]00NX}{[1][0][]source}{}{}{}}
\makeatother
\crefname{originallemma}{original lemma}{original lemmas}
\Crefname{originallemma}{Original lemma}{Original lemmas}
\makeatletter
\newlabel{algebra-lemma-flat-local-given-residue-field}{{03C3}{source}{}{sourcepointer-9}{}}
\newlabel{algebra-lemma-flat-local-given-residue-field@cref}{{[originallemma][0][]03C3}{[1][0][]source}{}{}{}}
\makeatother
\crefname{originallemma}{original lemma}{original lemmas}
\Crefname{originallemma}{Original lemma}{Original lemmas}
\makeatletter
\newlabel{lemma-proper-flat-h0}{{0E62}{source}{}{sourcepointer-10}{}}
\newlabel{lemma-proper-flat-h0@cref}{{[originallemma][0][]0E62}{[1][0][]source}{}{}{}}
\makeatother
\title{For a proper flat morphism with geometrically reduced fibre the pushforward splits off a finite etale algebra}
\author{The Stacks Project Authors}
\date{stacks-project:perfect.tex:local-snapshot-20261012; GNU FDL 1.2 or later}
\begin{document}
\maketitle
\setcounter{subsection}{1}
\noindent\textit{Editorial note.} All delivered mathematical text is the author's own, copied from the pinned Stacks Project chapter file perfect.tex; no statement or proof step was rewritten, reordered or removed. The counted span is the author environment and its own proof environment. References inside the span are handed to the assembler as navigation stubs that carry the canonical tag of the cited tag only; whether each of them is load-bearing is recorded on the stub itself.\par\medskip
\noindent\textit{Allowed standard prerequisites (editorial).}\begin{itemize}
\item Cohomology and base change makes the pushforward of the structure sheaf perfect with base change compatible formation, hence of tor amplitude bounded below by zero. invoked at proof lines 8414-8422. Provides the perfect complex whose decomposition is established.
\item After localizing at a suitable element a perfect complex splits off its degree zero cohomology with a perfect complement of tor amplitude bounded below by one. invoked at proof lines 8423-8435. Produces the direct sum decomposition in the statement.
\item The dimension of the degree zero cohomology of the fibre is semicontinuous, and equals the number of connected components of the geometrically reduced fibre. invoked at proof lines 8436-8470. Identifies the rank of the free module with the number of components and makes the condition on the fibre persist nearby.
\item An idempotent on the fibre comes from the pushforward of the structure sheaf, which gives the comparison between the algebra structure on the degree zero cohomology and the finite etale algebra of components. invoked at proof lines 8471-8480. Shows that the degree zero cohomology is the required finite etale algebra.
\end{itemize}\medskip
\ifdefined\hypertarget\hypertarget{sourcepointer-1}{}\fi
\noindent\textit{Source reference.} Original lemma 0B91: sources/stacks-perfect/pinned-original-source.tex, source lines 7742--7742. Invoked at proof lines 8414-8418 to see that the pushforward is perfect and commutes with base change; the decomposition argument built on that is inside the counted span.\par\medskip
\ifdefined\hypertarget\hypertarget{sourcepointer-2}{}\fi
\noindent\textit{Source reference.} Original lemma 02KH: sources/stacks-chow/tags.txt, source lines 1--1. Cited at proof lines 8419-8422 in support of the base change compatibility; the consequences drawn are inside the counted span.\par\medskip
\ifdefined\hypertarget\hypertarget{sourcepointer-3}{}\fi
\noindent\textit{Source reference.} Original lemma 0A1U: sources/stacks-chow/tags.txt, source lines 1--1. Used at proof lines 8429-8435 to obtain the local direct sum decomposition with the perfect complement; the computation of the rank is inside the counted span.\par\medskip
\ifdefined\hypertarget\hypertarget{sourcepointer-4}{}\fi
\noindent\textit{Source reference.} Original lemma 0658: sources/stacks-chow/tags.txt, source lines 1--1. Invoked at proof lines 8436-8440 to see that the degree zero cohomology is finite locally free; the rank computation that follows is inside the counted span.\par\medskip
\ifdefined\hypertarget\hypertarget{sourcepointer-5}{}\fi
\noindent\textit{Source reference.} Original lemma 0BUG: sources/stacks-chow/tags.txt, source lines 1--1. Cited at proof lines 8441-8445 to compute the rank of the free module on the fibre; the semicontinuity argument that keeps the rank nearby is inside the counted span.\par\medskip
\ifdefined\hypertarget\hypertarget{sourcepointer-6}{}\fi
\noindent\textit{Source reference.} Original lemma 0G7X: sources/stacks-perfect/pinned-original-source.tex, source lines 8334--8334. Invoked at proof lines 8466-8472 to extend the idempotents of the fibre to the pushforward, which is what makes the degree zero cohomology a finite etale algebra; the conclusion is inside the counted span.\par\medskip
\ifdefined\hypertarget\hypertarget{sourcepointer-7}{}\fi
\noindent\textit{Source reference.} Original lemma 0476: sources/stacks-chow/tags.txt, source lines 1--1. Cited at proof lines 8473-8478 to conclude that the algebra is finite etale over the base after shrinking; the identification of the algebra is inside the counted span.\par\medskip
\ifdefined\hypertarget\hypertarget{sourcepointer-8}{}\fi
\noindent\textit{Source reference.} Original lemma 00NX: sources/stacks-chow/tags.txt, source lines 1--1. A navigation pointer in the algebra identification; the identification itself is inside the counted span.\par\medskip
\ifdefined\hypertarget\hypertarget{sourcepointer-9}{}\fi
\noindent\textit{Source reference.} Original lemma 03C3: sources/stacks-chow/tags.txt, source lines 1--1. A navigation pointer used while comparing the algebra structures; the comparison is inside the counted span.\par\medskip
\ifdefined\hypertarget\hypertarget{sourcepointer-10}{}\fi
\noindent\textit{Source reference.} Original lemma 0E62: sources/stacks-perfect/pinned-original-source.tex, source lines 8264--8264. Cited in the opening sentence at proof lines 8414-8416 as reading advice for a similar proof; the argument for the present statement is written out in the counted span.\par\medskip
\begin{lemma}
\label{lemma-proper-flat-geom-red}
Let $f : X \to S$ be a morphism of schemes. Let $s \in S$. Assume
\begin{enumerate}
\item $f$ is proper, flat, and of finite presentation, and
\item the fibre $X_s$ is geometrically reduced.
\end{enumerate}
Then, after replacing $S$ by an open neighbourhood of $s$, there
exists a direct sum decomposition
$Rf_*\mathcal{O}_X = f_*\mathcal{O}_X \oplus P$
in $D(\mathcal{O}_S)$ where $f_*\mathcal{O}_X$ is a finite \'etale
$\mathcal{O}_S$-algebra and
$P$ is a perfect of tor amplitude in $[1, \infty)$.
\end{lemma}

\begin{proof}
The proof of this lemma is similar to the proof of
Lemma \ref{lemma-proper-flat-h0}
which we suggest the reader read first.
By cohomology and base change
(Lemma \ref{lemma-flat-proper-perfect-direct-image-general})
the complex $E = Rf_*\mathcal{O}_X$
is perfect and its formation commutes with arbitrary base change.
This first implies that $E$ has tor amplitude in $[0, \infty)$.

\medskip\noindent
We claim that after replacing $S$ by an open neighbourhood of
$s$ we can find a direct sum decomposition
$E = H^0(E) \oplus \tau_{\geq 1}E$ in $D(\mathcal{O}_S)$
with $\tau_{\geq 1}E$ of tor amplitude in $[1, \infty)$.
Assume the claim is true for now and assume we've made the replacement
so we have the direct sum decomposition.
Since $E$ has tor amplitude in $[0, \infty)$ we find that
$H^0(E)$ is a flat $\mathcal{O}_S$-module.
Hence $H^0(E)$ is a flat, perfect $\mathcal{O}_S$-module,
hence finite locally free, see
More on Algebra, Lemma \ref{more-algebra-lemma-perfect}
(and the fact that finite projective modules are finite locally free by
Algebra, Lemma \ref{algebra-lemma-finite-projective}).
Of course $H^0(E) = f_*\mathcal{O}_X$ is an $\mathcal{O}_S$-algebra.
By cohomology and base change we obtain
$H^0(E) \otimes \kappa(s) = H^0(X_s, \mathcal{O}_{X_s})$.
By Varieties, Lemma
\ref{varieties-lemma-proper-geometrically-reduced-global-sections}
and the assumption that $X_s$ is geometrically reduced, we
see that $\kappa(s) \to H^0(E) \otimes \kappa(s)$
is finite \'etale. By Morphisms, Lemma
\ref{morphisms-lemma-set-points-where-fibres-etale}
applied to the finite locally free morphism
$\underline{\Spec}_S(H^0(E)) \to S$,
we conclude that after shrinking $S$ the $\mathcal{O}_S$-algebra
$H^0(E)$ is finite \'etale.

\medskip\noindent
It remains to prove the claim. For this it suffices to prove that the map
$$
(f_*\mathcal{O}_X)_s
\longrightarrow
H^0(X_s, \mathcal{O}_{X_s}) = H^0(E \otimes^\mathbf{L} \kappa(s))
$$
is surjective, see
More on Algebra, Lemma \ref{more-algebra-lemma-better-cut-complex-in-two}.
Choose a flat local ring homomorphism $\mathcal{O}_{S, s} \to A$
such that the residue field $k$ of $A$ is algebraically closed, see
Algebra, Lemma \ref{algebra-lemma-flat-local-given-residue-field}.
By flat base change (Cohomology of Schemes, Lemma
\ref{coherent-lemma-flat-base-change-cohomology})
we get $H^0(X_A, \mathcal{O}_{X_A}) =
(f_*\mathcal{O}_X)_s \otimes_{\mathcal{O}_{S, s}} A$
and $H^0(X_k, \mathcal{O}_{X_k}) =
H^0(X_s, \mathcal{O}_{X_s}) \otimes_{\kappa(s)} k$.
Hence it suffices to prove that
$H^0(X_A, \mathcal{O}_{X_A}) \to H^0(X_k, \mathcal{O}_{X_k})$
is surjective. Since $X_k$ is a reduced proper scheme over $k$
and since $k$ is algebraically closed, we see that
$H^0(X_k, \mathcal{O}_{X_k})$ is a finite product of copies
of $k$ by the already used Varieties, Lemma
\ref{varieties-lemma-proper-geometrically-reduced-global-sections}.
Since by Lemma \ref{lemma-proper-idempotent-on-fibre}
the idempotents of this $k$-algebra are in the image
of $H^0(X_A, \mathcal{O}_{X_A}) \to H^0(X_k, \mathcal{O}_{X_k})$ we conclude.
\end{proof}

\end{document}
